% from main.tex
%%%%%%%%%%  scan idx 618  |  folio 605  %%%%%%%%%%
\[
  R^2f_*(I\otimes f^*M)\simeq \Hom_{\mathcal O_S}(Q,M).
\]
\item Le
conoyau \(E\) de \(f_*(\mathcal O_{X'})\to R^1f_*(I)\) est un \(\mathcal O_S\)-module localement libre.
\item L'homomorphisme \((f_*\mathcal O_X)_{\red}\longrightarrow (f_*\mathcal O_{X'})_{\red}\) est surjectif.
\end{enumerate}
Alors le morphisme canonique \(\Pic_{X/S}\longrightarrow \Pic_{X'/S}\) est représentable par un morphisme affine de présentation finie~; plus précisément, on a une suite exacte canonique de faisceaux f.p.p.f.,
\[
  0\longrightarrow \mathbb V(E^\vee)\longrightarrow \Pic_{X/S}
  \longrightarrow\Pic_{X'/S}\xrightarrow{u}\mathbb V(Q)
\]
et, si \(P=\Ker(u)\), le monomorphisme \(P\hookrightarrow\Pic_{X'/S}\) est représentable par une immersion fermée et \(\Pic_{X/S}\) est un torseur sur \(P\) sous le groupe \(\mathbb V(\check E)\) (c'est-à-dire, la suite de faisceaux f.p.p.f.
\[
  0\longrightarrow \mathbb V(\check E)\longrightarrow \Pic_{X/S}\longrightarrow P\longrightarrow0
\]
est exacte).

Considérons pour un \(S\)-schéma \(T\) les deux diagrammes suivants à lignes exactes :
\[
\begin{tikzcd}
0 \arrow[r] & I_T \arrow[r] & \mathcal O_{X_T} \arrow[r,dash] & \mathcal O_{X'_T} \arrow[r] & 0\\
0 \arrow[r] & I_T \arrow[r,"\alpha",dash] \arrow[u,equal] & \mathcal O_{X_T}^{*} \arrow[r] \arrow[u] & \mathcal O_{X'_T}^{*} \arrow[r] \arrow[u] & 0 ,
\end{tikzcd}
\qquad (\alpha(i)=1+i)
\]
et
\[
\begin{tikzcd}[column sep=tiny]
(f_T)_*(\mathcal O_{X'_T}) \arrow[r,"\delta_T"] &
R^1(f_T)_*(I_T) \arrow[d,equal]\\
(f_T)_*(\mathcal O_{X'_T}^{*}) \arrow[r,"\delta_T^*"] \arrow[u] &
R^1(f_T)_*(I_T) \arrow[r] &
R^1(f_T)_*(\mathcal O_{X_T}^{*}) \arrow[r] &
R'(f_T)_*(\mathcal O_{X'_T}^{*}) \arrow[r] &
R^2(f_T)_*(I_T) .
\end{tikzcd}
\]
On fera attention que le carré de gauche de ce diagramme \emph{n'est pas nécessairement commutatif}. Cependant on a~:
\[
  \textbf{3.2.}\qquad \im\delta_T^*=\im\delta_T.
\]
